\section{Evidence from Mechanistic Interpretability}
\label{sec:interp-evidence}

The mechanistic interpretability literature provides some of the
strongest empirical motivation for viewing transformers through a
message-passing lens. Importantly, most of this work was not originally
framed in terms of belief propagation or Bayesian networks. The value of
the BP perspective is therefore not that it replaces the
interpretability literature, but that it offers a common language for
organizing many of its findings.

This section surveys several influential interpretability results and
discusses how they naturally fit into a routing-and-update view of
transformer computation.

\subsection*{Induction Heads}

Olsson et al.~\cite{olsson2022induction} identify ``induction heads'':
two-head circuits that implement in-context pattern matching. Head~1
(the ``prefix'' head) attends from position $t$ back to the previous
occurrence of the current token. Head~2 (the ``induction'' head) then
attends to the token that followed that previous occurrence.

From the tutorial BP perspective, this resembles a two-stage gather
operation. One head retrieves a historical context; the second uses that
retrieved context to identify the next relevant source of evidence. The
important point is not that induction heads are ``literally'' BP nodes,
but that the routing structure closely resembles iterative message
passing over an implicit dependency graph.

\subsection*{Name Mover Heads and the IOI Circuit}

Wang et al.~\cite{wang2022ioi} reverse-engineer the indirect object
identification circuit in GPT-2 small. The key components are ``name
mover heads'' that copy token identity information from one position to
another. The Q/K circuits of these heads exhibit approximately rank-1
structure, while the V circuits route token-specific content into target
positions.

This behavior is naturally compatible with the routing interpretation
developed earlier in the paper. In the constructive BP mapping, sparse
projectDim/crossProject-style operations perform a similar function:
matching one feature dimension and routing information into another
location in the residual stream.

The BP framework therefore provides one possible interpretation of why
such sparse routing patterns repeatedly emerge in trained models.

\subsection*{FFN Key-Value Memories}

Geva et al.~\cite{geva2021ffn} show that FFN layers behave like
key-value memories: the first-layer weight matrix detects patterns and
the second-layer matrix writes structured updates back into the residual
stream. Many hidden units correspond to human-interpretable concepts or
contexts.

Within the tutorial framework developed here, FFNs can be interpreted as
performing local update operations conditioned on gathered evidence. The
``key'' determines when a particular update should activate; the
``value'' determines what modification is written into the evolving
residual-stream state.

This interpretation does not require treating FFN rows as literal
conditional-probability tables. Rather, it suggests that the key-value
memory view and the message-passing view may be complementary ways of
describing the same computational structure.

\subsection*{ROME and Localized Editing}

Meng et al.~\cite{meng2022rome} demonstrate that factual associations in
GPT models can sometimes be edited by modifying localized FFN
parameters. The resulting edits are often surprisingly surgical:
changing one association without broadly disrupting neighboring
behavior.

From the BP perspective, this resembles modifying a localized update
rule or factor potential within a larger inference system. Again, the
point is not that ROME directly proves the existence of explicit Bayes
net tables inside the transformer. The point is that localized,
structured update behavior is naturally compatible with a factorized
message-passing interpretation.

\subsection*{Sparse Autoencoders and Features}

Templeton et al.~\cite{templeton2024monosemanticity} train sparse
autoencoders on Claude 3 Sonnet and recover millions of interpretable
features from the residual stream. Many features exhibit coherent
semantic structure, feature neighborhoods, and suppression
relationships.

Within the tutorial framing, these features can be viewed as candidate
nodes or directions in an implicit computational graph. Their activation
values behave like evolving internal state variables whose interactions
are mediated through attention and FFN updates.

The ``Golden Gate Claude'' experiment is especially suggestive from this
perspective: clamping one feature propagates structured downstream
effects through the model, much like injecting evidence into a connected
message-passing system.

\subsection*{Attribution Graphs}

Recent work from Anthropic~\cite{anthropic2025attribution} constructs
attribution graphs tracing information flow between features across
layers. These graphs make visible the pathways through which activations
influence one another during a computation.

From the tutorial BP perspective, attribution graphs resemble partial
visualizations of the model's implicit message-passing structure:
attention routes information, FFNs apply transformations, and residual
streams accumulate evolving state over depth.

\subsection*{Synthesis}

The BP perspective should be understood as an organizing framework for
these findings rather than a claim that interpretability researchers
were unknowingly ``discovering Bayes nets.''

The important observation is simply that many empirical findings in
mechanistic interpretability fit naturally into a routing-plus-update
view of transformer computation:
\begin{itemize}
\item induction heads resemble multi-stage retrieval,
\item routing heads move structured information between positions,
\item FFNs apply localized updates,
\item sparse features behave like persistent latent variables,
\item and attribution graphs reveal structured pathways of influence.
\end{itemize}

The convergence between these findings and the message-passing
perspective is one reason the BP framework is useful as a conceptual
tool for understanding transformers.